\documentclass[11pt]{article}
\usepackage{amsmath,amssymb,bm}
\usepackage[margin=25mm]{geometry}
\title{V17C.22: Nonadiabatic extension of the M3 micromotion closure}
\author{}
\date{}
\begin{document}
\maketitle

\section{Scope and notation}
This note derives the terms omitted when the converged V17C.21.1 M3
Fourier--collocation solution is treated as a frozen function of the slow
orbit.  Time is $\tau=\Omega_1 t$.  The prescribed RF phases satisfy
\begin{equation}
 \dot\psi_1=1,\qquad \dot\psi_2=\nu,
\end{equation}
and the slow orbital phase of particle $i$ satisfies $\dot\theta_i=\omega_i$.
For a moving-frame representation write
\begin{equation}
 \mathbf r_{i,\perp}=e^{i\theta_i}\bigl(r_i+\xi_i(\theta_i,\psi_1,\psi_2;X_i)\bigr),
 \qquad z_i=Z_i+\zeta_i,
 \qquad X_i=(r_i,u_i,\omega_i,Z_i,w_i).
\end{equation}
The dependence on all slow coordinates is explicit.  The current M3 closure
sets $\dot X_i=0$ while solving its fast Fourier equation.

\section{Material derivative of the Fourier response}
Define the phase transport operator
\begin{equation}
 \mathcal L_i(X_i)=\omega_i\partial_{\theta_i}
                 +\partial_{\psi_1}+\nu\partial_{\psi_2}.
\end{equation}
Along a nonstationary slow orbit, the derivative is not $\mathcal L_i$ alone:
\begin{equation}
 \mathcal D_i=\mathcal L_i(X_i)+\dot X_i\cdot\nabla_{X_i}.
\end{equation}
Consequently,
\begin{align}
 \dot\xi_i&=\mathcal L_i\xi_i+\xi_{i,X}\dot X_i,\\
 \ddot\xi_i&=\mathcal L_i^2\xi_i
 +\dot\omega_i\,\partial_{\theta_i}\xi_i
 +2\mathcal L_i(\xi_{i,X}\dot X_i)
 +\xi_{i,X}\ddot X_i
 +\xi_{i,XX}[\dot X_i,\dot X_i].
 \label{eq:material-second}
\end{align}
The term $\dot\omega_i\partial_{\theta_i}\xi_i$ is the derivative of the
phase-transport operator itself.  The last two terms are the acceleration
and curvature terms of the slow-state-dependent response.  Formula
\eqref{eq:material-second} follows by applying $\mathcal D_i$ twice; it does
not assume a perturbation expansion.

\section{Moving-frame acceleration}
Let $J$ denote a $90^\circ$ rotation in the transverse plane, equivalently
multiplication by $i$.  Differentiating $e^{i\theta_i}(r_i+\xi_i)$ gives
\begin{align}
 e^{-i\theta_i}\ddot{\mathbf r}_{i,\perp}
 ={}&\ddot r_i-r_i\omega_i^2
 +i(2\omega_i u_i+r_i\dot\omega_i)\\
 &+\ddot\xi_i+2i\omega_i\dot\xi_i
 +(i\dot\omega_i-\omega_i^2)\xi_i.
 \label{eq:rot-accel}
\end{align}
The axial acceleration is $\ddot z_i=\ddot Z_i+\ddot\zeta_i$.  The terms
proportional to $\omega_i$, $\dot\omega_i$, $u_i$, and $\ddot r_i$ must be
kept consistently in a transient closure.  At a circular averaged root
($u_i=w_i=\dot X_i=0$), the nonadiabatic material-derivative terms vanish;
therefore this correction changes the linear response and relaxation, not
the defining stationary root to first order.

\section{Nonlinear drag response}
For the adopted Cunningham-corrected Stokes--Schiller--Naumann model,
\begin{equation}
 \mathbf d_i(\mathbf v)=-\lambda_i\left[1+0.15(\chi_i|\mathbf v|)^{0.687}\right]\mathbf v,
 \qquad \chi_i=\mathrm{Re}_{i,\mathrm{per\ speed}}.
\end{equation}
Its velocity derivative is
\begin{equation}
 D_{\mathbf v}\mathbf d_i=-\lambda_i\left[
 (1+0.15\chi_i^{0.687}|\mathbf v|^{0.687})I
 +0.10305\chi_i^{0.687}|\mathbf v|^{-1.313}
   \mathbf v\mathbf v^{\mathsf T}\right],
 \label{eq:drag-jac}
\end{equation}
with the continuous limiting value at $|\mathbf v|=0$.  The velocity used in
\eqref{eq:drag-jac} contains the slow orbital velocity, the rotating-frame
terms, and $\dot\xi_i=\mathcal L_i\xi_i+\xi_{i,X}\dot X_i$.  Thus a first
nonadiabatic correction changes both the mean drag force and mean drag torque;
replacing it by a scalar drag evaluated only at the mean velocity is not
consistent with Schiller--Naumann nonlinearity.

\section{Dynamic Fourier--Galerkin closure}
Let $\mathcal R_i(\xi_i;X_i,\psi_1,\psi_2)$ denote the original V21.1
discretized fast residual, including all eight complex MFS modes and the
Fourier projection of the nonlinear drag.  The frozen closure solves
$\mathcal R_i=0$ after selecting the fast subspace.  A dynamic M3+ closure
replaces the frozen derivative occurrences by \eqref{eq:material-second}
and \eqref{eq:rot-accel}, evaluates drag at the corresponding full velocity,
and solves jointly with the slow force/torque balances.  If
\begin{equation}
 \mathcal G_i(\xi_i,X_i)=0
\end{equation}
is differentiable and its fast-subspace Jacobian is invertible, then
\begin{equation}
 \xi_{i,X}=-\mathcal G_{i,\xi}^{-1}\mathcal G_{i,X}.
\end{equation}
Second derivatives follow by differentiating this equation again or by
second-order automatic differentiation.  No empirical fit coefficient is
introduced.  The coupled solve must include the unknown $\dot X_i$ and
$\ddot X_i$ consistently; evaluating the correction with arbitrary prescribed
accelerations is not a closed slow-flow model.

\section{Pair phases and Coulomb force}
The V21.1 two-particle closure uses a product average over independent ring
phases.  This is appropriate only when both phases are nonresonant and sample
their torus uniformly over the averaging window.  The full pair Coulomb force
is
\begin{equation}
 \mathbf F^C_1=C_{12}\frac{\mathbf r_1-\mathbf r_2}
 {|\mathbf r_1-\mathbf r_2|^3},\qquad
 \mathbf F^C_2=-\frac{m_1}{m_2}\mathbf F^C_1.
\end{equation}
For a phase-resolved M3+ computation, its Fourier projection must retain both
orbital angles $(\theta_1,\theta_2)$, or a resonant relative phase when a
combination frequency is slow.  An independent-ring Coulomb average cannot
represent phase locking, close-approach modulation, or the tangential torque
along such a locked pair orbit.

\section{Invariant-set formulation and stability}
For the full 12-dimensional mechanical model define the extended autonomous
system
\begin{equation}
 \dot x=F(x,\psi_1,\psi_2),\qquad \dot\psi_1=1,\quad\dot\psi_2=\nu.
\end{equation}
A candidate smooth invariant torus parameterization $K$ with phase vector
$\varphi$ and frequency vector $\varpi$ must solve
\begin{equation}
 DK(\varphi)\varpi=F_{\mathrm{ext}}(K(\varphi)).
 \label{eq:invariance}
\end{equation}
The phase vector can contain the two external RF angles and zero, one, or two
independent orbital angles.  Its dimension is a result to be tested, not
assumed.  A Fourier parameterization of \eqref{eq:invariance}, with gauge
conditions fixing phase origins, provides a direct route to an invariant
object.  The linear cocycle $\dot\eta=D_xF(K(\varphi),\psi)\eta$ is then
decomposed into tangent and normal components to measure normal attraction.
An external phase-neutral exponent is not a mechanical instability; conversely,
two small conditional mechanical exponents alone do not establish two torus
angles.

\section{Applicability test}
At the nominal root, the orbital rates are about $10^{-3}\,\tau^{-1}$,
while the RF phase rates are $1$ and $3.2197\,\tau^{-1}$.  This gives a
clear frequency gap for the stationary orbit.  However, the least damped
adiabatic slow-flow mode is $-0.10593\,\tau^{-1}$, whose e-fold time is only
$9.44\,\tau=1.50$ RF1 periods.  The V21.1 fast/slow mask excludes modes below
$\kappa_c=0.05\,\tau^{-1}$; therefore the ratio
$|\Re\lambda_{\mathrm{slow}}|/\kappa_c\simeq2.12$ is not small for the
boundary of the excluded band.  The force spectrum and response gain of each
near-resonant mode must be included before assigning an error bound.  The
available diagnostics do not yet supply such a bound or a converged M3+
solution.

\section{Conclusion}
The equations above define the required nonadiabatic extension and show
precisely why the old equilibrium Jacobian is insufficient for predicting
full-system transient recovery.  The present V17C.22 results validate the
full-system variation equations and quantify finite-time weak directions, but
do not solve the dynamic Fourier closure or the invariant equation
\eqref{eq:invariance}.  Accordingly, M3+ remains a derived, testable model
formulation rather than a numerically validated stability predictor.
\end{document}
